\documentclass[10pt,twocolumn]{article}
\usepackage[margin=0.48in]{geometry}
\usepackage{amsmath,amssymb,mathtools,bm}
\usepackage{graphicx}
\usepackage[most]{tcolorbox}
\usepackage{booktabs}
\usepackage{enumitem}
\usepackage{xcolor}
\usepackage{hyperref}
\usepackage{caption}
\usepackage{microtype}
\hypersetup{colorlinks=true,linkcolor=blue!45!black,urlcolor=blue!45!black,citecolor=blue!45!black}
\setlength{\columnsep}{0.18in}
\setlength{\parindent}{0pt}
\setlength{\parskip}{2.2pt}
\setlist[itemize]{leftmargin=1.1em,itemsep=1pt,topsep=2pt}
\setlist[enumerate]{leftmargin=1.3em,itemsep=1pt,topsep=2pt}
\captionsetup{font=footnotesize,skip=3pt}
\allowdisplaybreaks
\newcommand{\R}{\mathbb R}
\newcommand{\range}{\operatorname{range}}
\newcommand{\diag}{\operatorname{diag}}
\newcommand{\spanop}{\operatorname{span}}
\newcommand{\nov}{\mathcal N}
\newcommand{\eps}{\varepsilon}
\newcommand{\ket}[1]{\lvert #1\rangle}
\newcommand{\bra}[1]{\langle #1\rvert}
\newcommand{\braket}[2]{\langle #1\mid #2\rangle}
\definecolor{deepblue}{HTML}{1E5575}
\definecolor{pale}{HTML}{EFF7FA}
\definecolor{warm}{HTML}{FFF4E8}
\definecolor{greenbox}{HTML}{EEF8EF}
\definecolor{redbox}{HTML}{FFF0F0}
\tcbset{boxsep=2pt,left=4pt,right=4pt,top=3pt,bottom=3pt,arc=2pt,boxrule=0.55pt}
\newtcolorbox{intuitionbox}{colback=pale,colframe=deepblue,title=Intuition,fonttitle=\bfseries\small}
\newtcolorbox{jargonbox}{colback=warm,colframe=orange!60!black,title=Jargon-heavy statement,fonttitle=\bfseries\small}
\newtcolorbox{cautionbox}{colback=redbox,colframe=red!55!black,title=Caution,fonttitle=\bfseries\small}

\title{Paper-III QSE Projection Geometry Primer\\[-1mm]
\large Same Gram/Schur machinery, written only in QSE notation}
\author{Local pedagogical note}
\date{May 17, 2026}
\begin{document}
\maketitle

\begin{abstract}
This note rewrites the nonorthogonal projection primer entirely in Paper-III QSE notation. The recurring object is a selected excitation basis \(\Phi_{\mathcal A_\ell}\), a candidate excitation vector \(\phi_r\), the QSE overlap matrix \(S_\ell\), and the candidate overlap vector \(s_r\). QSE novelty is the quotient residual
\[
F_r-s_r^\top S_{\ell,\tau}^+s_r,
\]
while Phase-III spectral scoring compares the residualized candidate with the outside-span Ritz residual of a target root.
\end{abstract}

\section{The QSE projection question}

\begin{intuitionbox}
The selected QSE basis already spans a nonorthogonal excitation subspace. A candidate excitation record is useful when its vector \(\phi_r\) has a component that survives projection against the selected basis and when that surviving component helps the target Ritz roots.
\end{intuitionbox}

\begin{jargonbox}
Paper III is Rayleigh--Ritz/Galerkin projection in a nonorthogonal excitation basis. The matrix \(S_\ell\) is the metric of the selected excitation subspace, \(M_\ell\) is the Hamiltonian projected into that subspace, and \(S_\ell^{-1}\)-type solves convert overlap data into basis coefficients. Candidate selection uses the same quotient residual:
\[
\|\bar\phi_r\|^2=F_r-s_r^\top S_{\ell,\tau}^+s_r.
\]
\end{jargonbox}

At QSE construction step \(\ell\), let
\begin{equation}
\Phi_{\mathcal A_\ell}=[\phi_a]_{a\in\mathcal A_\ell},
\qquad
\phi_r=Q_{\mathfrak s,0}\Omega_r\ket{\psi_0}.
\label{eq:p3_basis_candidate}
\end{equation}
The selected basis explains \(\phi_r\) by choosing coefficients \(\beta\):
\begin{equation}
\Phi_{\mathcal A_\ell}\beta=\sum_{a\in\mathcal A_\ell}\beta_a\phi_a\approx\phi_r.
\label{eq:p3_synthesis}
\end{equation}
Choose \(\beta\) by least squares:
\begin{equation}
\beta_r^\star\in\arg\min_\beta
\|\phi_r-\Phi_{\mathcal A_\ell}\beta\|^2.
\label{eq:p3_ls}
\end{equation}
Expanding gives
\begin{align}
\|\phi_r-\Phi_{\mathcal A_\ell}\beta\|^2
&=\langle\phi_r-\Phi_{\mathcal A_\ell}\beta,\,
\phi_r-\Phi_{\mathcal A_\ell}\beta\rangle
\notag\\
&=F_r-2\beta^\top s_r+\beta^\top S_\ell\beta,
\label{eq:p3_expand}\\
(S_\ell)_{ab}&=\phi_a^\top\phi_b,
\qquad
(s_r)_a=\phi_a^\top\phi_r,
\qquad
F_r=\phi_r^\top\phi_r.
\label{eq:p3_SsF}
\end{align}
The vector \(s_r\) records the candidate's overlaps against selected excitation vectors; it is analysis data. The solve through \(S_\ell\) converts those overlaps into coefficients in the nonorthogonal selected basis.
Stationarity gives
\begin{equation}
S_\ell\beta_r^\star=s_r,
\qquad
\beta_r^\star=S_{\ell,\tau}^+s_r.
\label{eq:p3_normal}
\end{equation}
The metric projector is
\begin{equation}
P_{\ell,\tau}=\Phi_{\mathcal A_\ell}S_{\ell,\tau}^+\Phi_{\mathcal A_\ell}^\top.
\label{eq:p3_projector}
\end{equation}
The residualized candidate is
\begin{equation}
\bar\phi_r=(I-P_{\ell,\tau})\phi_r,
\qquad
\|\bar\phi_r\|^2=F_r-s_r^\top S_{\ell,\tau}^+s_r.
\label{eq:p3_residualized_candidate}
\end{equation}
The normalized QSE novelty is
\begin{equation}
\boxed{
\nov_{\rm QSE}(r;\ell)=1-\frac{s_r^\top S_{\ell,\tau}^+s_r}{F_r}.
}
\label{eq:p3_novelty}
\end{equation}

\section{Schur and quotient readings in QSE notation}

Append the candidate vector to the selected excitation basis. The augmented overlap matrix is
\begin{equation}
\mathcal S_{\ell r}
=
\begin{pmatrix}
\Phi_{\mathcal A_\ell}^\top\Phi_{\mathcal A_\ell}
&\Phi_{\mathcal A_\ell}^\top\phi_r\\
\phi_r^\top\Phi_{\mathcal A_\ell}
&\phi_r^\top\phi_r
\end{pmatrix}
=
\begin{pmatrix}
S_\ell&s_r\\
s_r^\top&F_r
\end{pmatrix}.
\label{eq:p3_augmented_overlap}
\end{equation}
When \(S_\ell\) is invertible, the candidate's Schur residual is
\begin{equation}
\mathcal S_{\ell r}/S_\ell
=F_r-s_r^\top S_\ell^{-1}s_r.
\label{eq:p3_schur_residual}
\end{equation}
With the supported inverse, the quotient-space version is
\begin{equation}
\begin{aligned}
\phi_r\sim\phi_r'
&\Longleftrightarrow
\phi_r-\phi_r'\in\range(\Phi_{\mathcal A_\ell}),\\
\|[\phi_r]\|^2
&:=\min_\beta\|\phi_r-\Phi_{\mathcal A_\ell}\beta\|^2\\
&=F_r-s_r^\top S_{\ell,\tau}^+s_r.
\end{aligned}
\label{eq:p3_quotient_residual}
\end{equation}
Thus QSE novelty asks how much of the candidate excitation vector remains after selected basis vectors have had their best coefficient-space explanation.

\section{Supported QSE metric}

Diagonalize the selected-basis overlap matrix:
\begin{equation}
S_\ell=V_\ell\diag(\sigma_{\ell,1},\ldots,\sigma_{\ell,m})V_\ell^\top.
\label{eq:p3_eigs}
\end{equation}
The columns \(v_{\ell,a}\) are obtained from
\begin{equation}
S_\ell v_{\ell,a}=\sigma_{\ell,a}v_{\ell,a},
\qquad \|v_{\ell,a}\|=1.
\label{eq:p3_eigenproblem}
\end{equation}
Each \(v_{\ell,a}\) is a coefficient recipe for a selected-basis mixture:
\begin{equation}
\Phi_{\mathcal A_\ell}v_{\ell,a}=\sum_{b\in\mathcal A_\ell}(v_{\ell,a})_b\phi_b.
\label{eq:p3_mode_synthesis}
\end{equation}
Its physical norm is
\begin{equation}
\|\Phi_{\mathcal A_\ell}v_{\ell,a}\|^2
=v_{\ell,a}^\top S_\ell v_{\ell,a}=\sigma_{\ell,a}.
\label{eq:p3_mode_length}
\end{equation}
The supported inverse is
\begin{equation}
S_{\ell,\tau}^+
=V_\ell\diag\left(\frac{\mathbf 1[\sigma_{\ell,a}\ge\tau\sigma_{\ell,\max}]}{\sigma_{\ell,a}}\right)V_\ell^\top.
\label{eq:p3_supported_inverse}
\end{equation}

\section{QSE roots are coefficient mixtures}

Define the projected Hamiltonian and overlap matrices:
\begin{equation}
M_\ell=\Phi_{\mathcal A_\ell}^\top H\Phi_{\mathcal A_\ell},
\qquad
S_\ell=\Phi_{\mathcal A_\ell}^\top\Phi_{\mathcal A_\ell}.
\label{eq:p3_MS}
\end{equation}
The QSE root \(\nu\) is synthesized as
\begin{equation}
\Psi_\nu^{(\ell)}=\Phi_{\mathcal A_\ell}c^{(\nu)}
=\sum_{a\in\mathcal A_\ell}c_a^{(\nu)}\phi_a.
\label{eq:p3_root}
\end{equation}
The coefficients solve
\begin{equation}
\boxed{M_\ell c^{(\nu)}=\epsilon_\nu S_\ell c^{(\nu)}.}
\label{eq:p3_gep}
\end{equation}
The vector \(c^{(\nu)}\) is the coefficient recipe for the approximate root inside the selected nonorthogonal basis.

The generalized eigenproblem comes from stationarity of the Rayleigh quotient
\begin{equation}
\epsilon(c)=\frac{c^\top M_\ell c}{c^\top S_\ell c}.
\label{eq:p3_rayleigh}
\end{equation}
Differentiating under the normalization \(c^\top S_\ell c=1\) gives
\begin{equation}
\delta\left(c^\top M_\ell c-\epsilon\,c^\top S_\ell c\right)=0
\quad\Longrightarrow\quad
M_\ell c=\epsilon S_\ell c.
\label{eq:p3_gep_derivation}
\end{equation}
The metric \(S_\ell\) appears because the basis vectors are nonorthogonal; coefficient dot products must be measured with the overlap metric.

\section{Outside-span residual}

The full eigen-equation error for root \(\nu\) is
\begin{equation}
y_\nu^{(\ell)}=(H-\epsilon_\nu)\Psi_\nu^{(\ell)}.
\label{eq:p3_y}
\end{equation}
Project it outside the selected QSE span:
\begin{equation}
R_\nu^{(\ell)}=(I-P_{\ell,\tau})y_\nu^{(\ell)}
=(I-P_{\ell,\tau})(H-\epsilon_\nu)\Phi_{\mathcal A_\ell}c^{(\nu)}.
\label{eq:p3_residual}
\end{equation}
This is the outside-span part of the eigen-equation error. If \(H\Psi_\nu^{(\ell)}\) stays inside the selected QSE span and equals \(\epsilon_\nu\Psi_\nu^{(\ell)}\) there, then this residual is small.
Residual capture compares \(\bar\phi_r\) to \(R_\nu^{(\ell)}\):
\begin{equation}
C_\nu^{\rm res}(r;\ell)
=
\frac{|\bar\phi_r^\top R_\nu^{(\ell)}|^2}
{(\bar\phi_r^\top\bar\phi_r+\lambda_S)((R_\nu^{(\ell)})^\top R_\nu^{(\ell)}+\eps_R)}.
\label{eq:p3_res_capture}
\end{equation}

\section{Two-dimensional Ritz gain}

Normalize the surviving candidate direction:
\begin{equation}
u_r=\frac{\bar\phi_r}{\sqrt{\bar\phi_r^\top\bar\phi_r+\lambda_S}}.
\label{eq:p3_ur}
\end{equation}
Define the root-candidate coupling and shifted separation:
\begin{equation}
\eta_{\nu r}=u_r^\top(H-\epsilon_\nu)\Psi_\nu^{(\ell)},
\qquad
\delta_{\nu r}=u_r^\top(H-\epsilon_\nu)u_r.
\label{eq:p3_eta_delta}
\end{equation}
\(\eta_{\nu r}\) is the mixing between the current Ritz root and the residualized candidate direction. \(\delta_{\nu r}\) is the candidate's shifted Rayleigh energy relative to the current root.
The shifted Hamiltonian in \(\spanop\{\Psi_\nu^{(\ell)},u_r\}\) is
\begin{equation}
\begin{pmatrix}
0&\eta_{\nu r}\\
\eta_{\nu r}^*&\delta_{\nu r}
\end{pmatrix}.
\label{eq:p3_2x2_block}
\end{equation}
The lower-branch improvement is
\begin{equation}
\boxed{
G_{\nu r}^{2\times2}
=\frac{\sqrt{\delta_{\nu r}^2+4|\eta_{\nu r}|^2}-\delta_{\nu r}}{2}
\approx\frac{|\eta_{\nu r}|^2}{\delta_{\nu r}}.
}
\label{eq:p3_2x2_gain}
\end{equation}
Thus Phase III asks whether the candidate direction that survived basis projection also couples to the target root strongly enough to improve its Ritz energy.

\appendix
\section{Tiny QSE metric example}

Let two selected excitation vectors be nearly equal:
\begin{equation}
\phi_1=\begin{pmatrix}1\\0\end{pmatrix},
\qquad
\phi_2=\begin{pmatrix}1\\\delta\end{pmatrix},
\qquad 0<\delta\ll1.
\end{equation}
Then
\begin{equation}
S=\begin{pmatrix}1&1\\1&1+\delta^2\end{pmatrix},
\qquad
v_-\approx\frac1{\sqrt2}\begin{pmatrix}1\\-1\end{pmatrix},
\qquad
\sigma_-\approx\frac{\delta^2}{2}.
\end{equation}
The difference recipe \(v_-\) synthesizes a tiny physical vector, so support truncation removes it before solving the GEP.

\section{Example: ordinary eigenproblem before QSE}

For
\begin{equation}
A=
\begin{pmatrix}
2&0\\0&5
\end{pmatrix},
\qquad
Av=\lambda v,
\qquad
v=\begin{pmatrix}x\\y\end{pmatrix},
\end{equation}
the equations are
\begin{equation}
2x=\lambda x,\qquad 5y=\lambda y.
\end{equation}
Thus the two roots are
\begin{equation}
\lambda_0=2,\quad v_0=(1,0)^\top,
\qquad
\lambda_1=5,\quad v_1=(0,1)^\top.
\end{equation}
QSE replaces \(A\) by the matrix pencil \((M_\ell,S_\ell)\), and the eigenvector \(c^{(\nu)}\) is the coefficient recipe for the approximate quantum state \(\Psi_\nu^{(\ell)}\).

\section{Example: two-vector generalized eigenproblem}

Let
\begin{equation}
S=
\begin{pmatrix}
1&1/2\\
1/2&1
\end{pmatrix},
\qquad
M=
\begin{pmatrix}
1&0\\
0&3
\end{pmatrix}.
\end{equation}
The generalized eigenproblem is
\begin{equation}
Mc=\epsilon Sc.
\end{equation}
Equivalently,
\begin{equation}
\det(M-\epsilon S)=0.
\end{equation}
Here
\begin{equation}
\det
\begin{pmatrix}
1-\epsilon&-\epsilon/2\\
-\epsilon/2&3-\epsilon
\end{pmatrix}
=(1-\epsilon)(3-\epsilon)-\frac{\epsilon^2}{4}=0.
\end{equation}
The overlap metric changes the root equation because coefficient vectors are measured in a nonorthogonal basis.

\section{Example: residual capture}

Let a normalized residualized candidate and outside-span residual satisfy
\begin{equation}
\bar\phi_r^\top\bar\phi_r=1,\qquad
(R_\nu^{(\ell)})^\top R_\nu^{(\ell)}=4,\qquad
\bar\phi_r^\top R_\nu^{(\ell)}=1.
\end{equation}
Ignoring small stabilizers, the residual-capture score is
\begin{equation}
C_\nu^{\rm res}(r;\ell)=\frac{1^2}{1\cdot4}=\frac14.
\end{equation}
The candidate captures one quarter of the residual direction's squared normalized alignment.

\section{Example: two-dimensional Ritz gain}

Let
\begin{equation}
\eta_{\nu r}=0.3,\qquad
\delta_{\nu r}=3.
\end{equation}
Then
\begin{equation}
G_{\nu r}^{2\times2}
=\frac{\sqrt{3^2+4(0.3)^2}-3}{2}
=\frac{\sqrt{9.36}-3}{2}
\approx 0.0297.
\end{equation}
The perturbative estimate gives
\begin{equation}
\frac{|\eta_{\nu r}|^2}{\delta_{\nu r}}
=\frac{0.09}{3}=0.03.
\end{equation}
The candidate produces a lower-branch Ritz improvement of about \(3\times10^{-2}\) in the same energy units.

\end{document}
